% === IX. Dokumentasi Cara Penggunaan Perangkat Lunak ===
\secbox{IX.}{Dokumentasi Cara Penggunaan Perangkat Lunak}

Secara umum, penggunaan LUNA dilakukan melalui tahapan berikut:
\begin{enumerate}[leftmargin=*,label=\arabic*.,itemsep=0.15em,parsep=0pt]
    \item Pengguna melakukan masuk (\textit{login}) ke dalam aplikasi.
    \item Pengguna memulai percakapan menggunakan teks atau suara.
    \item Sistem menganalisis percakapan melalui Multi-Agent Artificial Intelligence untuk memahami konteks dan kondisi emosional pengguna.
    \item LUNA memberikan respons pendampingan yang dipersonalisasi sesuai hasil analisis.
    \item Hasil percakapan secara otomatis disusun menjadi diary harian dan disimpan sebagai riwayat pengguna.
    \item Sistem melakukan analisis kondisi mental secara berkelanjutan serta memberikan rekomendasi yang sesuai dengan kebutuhan pengguna.
    \item Apabila terdeteksi indikasi risiko krisis mental, sistem menjalankan mekanisme \textit{Early Mental Crisis Detection} sesuai kebijakan yang diterapkan.
\end{enumerate}

Dokumentasi penggunaan secara rinci setiap langkah disajikan pada Lampiran 5.

\vspace{0.8em}
\subsecbar{1. Rencana Uji Coba Klinis \& Pilot Testing}

\indent Dalam rangka memastikan efektivitas serta keamanan respons pendampingan LUNA pada skenario dunia nyata, tim pengembang telah menginisiasi komunikasi dan kolaborasi awal bersama pakar kesehatan mental, \textbf{Qori Fanani, M.Psi., Psikolog} (Petugas Poliklinik Politeknik Negeri Malang dan Dosen Universitas Kepanjen). Pihak Poliklinik Polinema sangat terbuka untuk mendukung pelaksanaan \textit{clinical pilot testing} dan uji coba penggunaan LUNA kepada mahasiswa maupun pasien di poliklinik. Tahapan rencana uji coba klinis lanjutan yang dirancang meliputi:

\begin{enumerate}[leftmargin=*,label=\arabic*.,itemsep=0.15em,parsep=0pt]
    \item \textbf{Uji Validasi Respons AI oleh Psikolog Klinis:} Evaluasi tingkat empati, kesesuaian modul intervensi CBT, serta presisi kalkulasi skor DASS-21 (Level 1--3) oleh ahli kesehatan mental profesional.
    \item \textbf{Pilot Deployment Terbatas:} Penerapan aplikasi LUNA secara bertahap pada responden pengguna di Poliklinik Polinema di bawah pengawasan psikolog.
    \item \textbf{Evaluasi Dampak \& Refinement:} Pengukuran perubahan skor depresi, kecemasan, dan stres sebelum serta sesudah penggunaan LUNA untuk melakukan penyempurnaan alur agen AI secara berkelanjutan.
\end{enumerate}



